\section{Discussion}

\subsection{What the Certificate Establishes}

Dynamic capture exposes a persistence asymmetry.
A tracking error affects the current pose and may propagate, whereas an unsafe Gaussian write can remain visible from many future viewpoints.
Hard exclusion minimizes the second risk but may increase the first when a dynamic mask removes most image support.
GDOR asks whether an excluded observation is sufficiently consistent to support transient localization; it does not declare every recovered pixel truly static.

The non-vacuous route certificate sharpens this claim.
For Full, recovery occurs in every tested sequence-seed cell, a mapping weight exists on every recovered frame, and the exact recovered-support/map-admission intersection is empty.
The Full-NoM counterfactual then exposes overlap in every cell when the conservative route is removed.
This directly verifies the implemented privilege boundary.
It does not certify the semantic correctness of every admitted static pixel, map completeness, or zero 3D ghost contamination.

\subsection{Why the Ordered Ladder Is Non-Monotonic}

Full improves aggregate tracking over Semantic, but it is not the lowest-ATE row in the mechanism matrix.
The simpler T+M configuration recovers more support without the later flow veto, risk gate, or adaptive trigger and performs best on the consumed ten-sequence aggregate.
Adding guards can suppress useful support as well as unsafe support; their net effect depends on occlusion geometry, prediction quality, and tracker state.
The ladder therefore diagnoses complete operating points, not independent factor effects.

This finding does not invalidate Full's routing claim.
Both Full and T+M retain the conservative mapping route and pass every non-vacuous overlap certificate.
It instead narrows the next design question: whether T+M's simpler tracking policy generalizes beyond the sequences that revealed it.
Promoting T+M without a fresh split would convert development evidence into a held-out claim.

\subsection{Sequence and Mapping Boundaries}

The main result is driven by genuine rescues on difficult sequences, including TUM sitting\_halfsphere and Bonn crowd2, while Full loses on three of ten sequence means against Semantic.
MapMatched also beats Full on six sequence means despite a worse overall mean and failure rate.
These outcomes indicate that dynamic-region recovery should be described as regime-dependent support restoration rather than a universal replacement for conservative masking.

The mapping study is stronger than the earlier single-seed diagnostic because it uses four scenes, three seeds, common held-out views, and both online and GT-aligned rendering.
Its static-region PSNR and SSIM gains show that the tracking intervention does not require degrading the tested reconstructions.
Still, the scene count is limited, LPIPS is absent, and the occluded-background fields are partial proxies.
The appropriate claim is positive multi-seed common-view reconstruction evidence, not general photorealistic or ghost-free superiority.

\subsection{Reproducibility and External Validity}

The public evidence package contains path-sanitized sequence means, all 210 route-certificate cells, all 72 mapping cells, source/release hashes, and deterministic validators.
Raw datasets, trajectories, rendered frames, and Gaussian maps remain in the retained experiment release rather than the source repository.
The frozen August manuscript is preserved separately; this DYN-19 revision has its own evidence manifest and PDF identity.

External validity remains the principal ranking limitation.
Concurrent systems overlap with the high-level tracking--mapping separation premise, and no official DyPho-SLAM runner has been evaluated under the exact DYN-19 contract.
GDOR's defensible contribution is the concrete RGB-D/ORB routing mechanism, exact non-vacuous route audit, and same-source evidence.
Broader SOTA language requires a protocol-matched external baseline rather than comparison to published numbers alone.
